% =====================================================================
%  fr/blocks.tex — tous les blocs du CV, rédigés une seule fois, en français.
%  Les fichiers de variante (cv-*-fr.tex) choisissent les blocs et leur ordre.
%  Dans un bloc :
%     \HW{...}  n'apparaît que dans le CV « Ingénieur logiciel et matériel embarqué »
%     \SW{...}  n'apparaît que dans le CV « Ingénieur logiciel embarqué »
%  Les termes techniques anglais d'usage courant (thread, gateway, buck…) sont conservés.
% =====================================================================

% ---------------------------------------------------------------- en-tête
\newcommand{\BlkHeader}[1]{%
  \cvheader{Junior Deogracias MAMBO MOTSOU}{#1}
  {\href{mailto:motsoujunior17@gmail.com}{motsoujunior17@gmail.com}\sep
   \phone\sep France\sep \availability}
  {\href{https://teo1747.github.io/Portfolio/}{teo1747.github.io/Portfolio}\sep
   \href{https://github.com/teo1747}{github.com/teo1747}\sep
   \href{https://www.linkedin.com/in/teo1747}{linkedin.com/in/teo1747}}}
\newcommand{\phone}{+33 6 01 16 86 06}
\newcommand{\availability}{Disponible immédiatement}

% ---------------------------------------------------------------- profil
\newcommand{\BlkProfile}{%
\section{Profil}
\FP{Ingénieur systèmes embarqués (Polytech Dijon, 2026) orienté conception numérique sur FPGA et SoC Zynq :
RTL en VHDL, interfaces AXI4-Lite et AXI4-Stream, intégration d'IP et block design sous Vivado, et le
logiciel bare-metal qui pilote le matériel depuis Vitis. Je vérifie le RTL contre des modèles C bit-exacts
avec des bancs de test auto-vérifiants. Une année de mise en route firmware automobile chez FORVIA
m'apporte la vision système.}%
\NF{\HW{Ingénieur systèmes embarqués (Polytech Dijon, 2026) couvrant firmware, électronique et matériel numérique. Chez FORVIA, j'ai mené la mise en route du firmware sur de nouveaux composants
automobiles (contrôleur de chauffe LIN, nœuds LED MeLiBu) et développé le logiciel d'un démonstrateur
d'éclairage à neuf gateways. Je travaille à partir des datasheets, je localise les défauts côté matériel
ou logiciel et je valide sur banc ; la conception VHDL/FPGA complète ce profil.}%
\SW{Ingénieur logiciel embarqué (Polytech Dijon, 2026) spécialisé en firmware C, systèmes d'exploitation
temps réel et outillage associé. Chez FORVIA, j'ai développé et débogué du firmware pour
microcontrôleurs automobiles ainsi qu'une chaîne de pilotage multithread Python/Kotlin pour un démonstrateur
d'éclairage. Dans mes projets personnels, je développe un RTOS, un OS x86-64 et une chaîne de compilation C,
avec un soin particulier pour le build, la CI et les tests.}}\par}

% ---------------------------------------------------------------- compétences
\newcommand{\BlkSkillsHW}{%
\section{Compétences techniques}
\cvskill{Microcontrôleurs}{STM32 (F446RE, G474RE), Infineon TLD4020 (Cortex-M23), Melexis MLX81116/MLX81143,
  RP2040 (PIO), PIC32MZ, PIC24, ATmega328P}
\cvskill{Firmware}{C (bare-metal, registres via CMSIS, HAL/SDK), FreeRTOS, Zephyr, interruptions et timers,
  acquisition ADC, PWM, machines d'états}
\cvskill{Bus \& protocoles}{LIN, MeLiBu, UART, SPI, I²C, AXI4-Stream, Bluetooth série, ISO 14443-A (MIFARE)}
\cvskill{FPGA \& VHDL}{VHDL (machines d'états, chemins de données pipelinés, IOBUF bidirectionnel),
  Vivado (synthèse, implémentation, IP integrator, ILA), Vitis HLS, XSim ; Artix-7, Zynq-7000}
\cvskill{Électronique}{Schéma et routage PCB (KiCad, Proteus), interfaçage capteurs (CTN, capteurs I²C/SPI),
  mise en route de cartes d'évaluation}
\cvskill{Debug \& outils}{J-Link, ST-LINK, débogueur PICkit embarqué, GDB, consoles série ; IAR EWARM,
  MPLAB X + Harmony v3/MCC, STM32CubeProgrammer, Pico SDK ; Git, CMake, Ninja, West}
\cvskill{Logiciel}{Python (PyQt5, OpenCV, threading), C++/Qt, Kotlin (Android, Jetpack Compose)}\par}

\newcommand{\BlkSkillsSW}{%
\section{Compétences techniques}
\cvskill{Langages}{C (C99/C11), Python, C++ (Qt), Kotlin ; assembleur : x86-64 (NASM), Arm Thumb, RISC-V}
\cvskill{RTOS \& OS}{FreeRTOS (tâches, queues, mutex, event groups, stream buffers), Zephyr (devicetree,
  Kconfig), bare-metal ; noyau : ordonnancement, commutation de contexte, pagination, interruptions, syscalls}
\cvskill{Linux}{Ubuntu au quotidien ; drivers en espace utilisateur (spidev, i2c-dev, GPIO sysfs) sur Raspberry Pi}
\cvskill{Build \& toolchains}{GCC (arm-none-eabi, cross x86\_64-elf), Clang, scripts d'édition de liens,
  CMake, Ninja, GNU Make, West, Kconfig ; IAR EWARM, MPLAB X (XC16/XC32)}
\cvskill{Debug \& tests}{GDB (distant, QEMU), J-Link, QEMU ; tests unitaires sous ASan/UBSan, tests de
  conformité et tests différentiels contre des modèles de référence}
\cvskill{Automatisation}{Git (workflow par pull requests), CI GitHub Actions, outillage Python}
\cvskill{Interfaces}{UART, SPI, I²C, LIN, MeLiBu, Bluetooth série avec protocole JSON, clients REST/HTTP}
\cvskill{Applications}{Python multithread (PyQt5, OpenCV), Qt/C++, Android (Kotlin, Jetpack Compose)}\par}

\newcommand{\BlkSkillsFPGA}{%
\section{Compétences techniques}
\cvskill{HDL \& conception}{VHDL (RTL synthétisable, bancs de test VHDL-2008), Verilog (bases) ; machines
  d'états, chemins de données pipelinés, line buffers, gestion de la croissance en bits, handshakes AXI avec
  back-pressure, E/S bidirectionnelles}
\cvskill{Flot AMD/Xilinx}{Vivado : block design sous IP integrator, packaging d'IP, synthèse, implémentation,
  XDC, rapports de timing et d'utilisation, bitstream, ILA et Hardware Manager ; Vitis HLS ; applications
  bare-metal Vitis ; scripts TCL du flot}
\cvskill{Composants \& cartes}{Zynq-7000 (ZedBoard, xc7z020), Artix-7 (Nexys 4 DDR, xc7a100t)}
\cvskill{Interfaces \& IP}{AXI4-Lite, AXI4-Stream, AXI DMA, PS7 Zynq (ports GP/HP), Video Timing Controller,
  Clocking Wizard, FIFO Generator ; VGA, I²C, UART, SPI, SRAM ZBT}
\cvskill{Vérification}{Bancs de test auto-vérifiants, modèles C de référence bit-exacts, injection de
  blocages aléatoires ; GHDL, XSim, ModelSim ; simulation en CI sur GitHub Actions}
\cvskill{Logiciel embarqué}{C (bare-metal, drivers standalone Xilinx, CMSIS), FreeRTOS, Zephyr ; Python, TCL}
\cvskill{Debug}{ILA, JTAG / Hardware Manager, J-Link, GDB, consoles série}\par}

% ---------------------------------------------------------------- FORVIA
\newcommand{\BlkForvia}{%
\cventry{Ingénieur systèmes embarqués --- contrat de professionnalisation}
  {FORVIA -- Faurecia Interior Systems}{département Surface Activation, France}{\forviadates}
\HW{\BlkForviaTCU\BlkForviaMeli\BlkForviaLight}%
\SW{\BlkForviaLight\BlkForviaTCU\BlkForviaMeli}}

% version condensée (inutilisée pour l'instant ; à utiliser pour gagner de la place)
\newcommand{\BlkForviaShort}{%
\begin{itemize}
  \item Mise en route firmware sur Infineon TLD4020-3ET (Cortex-M23, LIN) pour un contrôleur d'élément
    chauffant : commandes LIN, PWM via le LDRV, régulation CTN ; HardFault attribué à une révision de
    silicium, échantillonnage ADC synchronisé sur la PWM (ISR à mi-impulsion, moyenne sur 16 mesures).
  \item Portage d'un driver LED MeLiBu sur RP2040 (machines d'états PIO) du MLX81116 vers le MLX81143 pour
    un panneau de 154 LED ; correction du timing bus (break, section critique, vidage de la FIFO TX).
  \item Responsable logiciel d'un démonstrateur d'éclairage intérieur : Python multithread pilotant neuf
    gateways UART STM32 synchronisées à 921\,600 bauds, pilotage par tablette Android en Bluetooth.
\end{itemize}
\cvenv{Environnement\,:}{IAR EWARM, J-Link, Pico SDK, Python/PyQt5, Kotlin, STM32, LIN, MeLiBu, UART}}
\newcommand{\forviadates}{sept.\ 2025 -- sept.\ 2026}

\newcommand{\BlkForviaTCU}{%
\cvsub{Mise en route firmware : contrôleur d'élément chauffant piloté en LIN (TCU)}
\begin{itemize}
  \item Développement du firmware C sur Infineon TLD4020-3ET (Arm Cortex-M23 avec LIN intégré) :
    commandes ON/OFF par LIN, PWM de chauffe via le périphérique LDRV, régulation de température par CTN,
    protections, architecture en machine d'états.
  \item Diagnostic des blocages de mise en route : HardFault à l'initialisation causé par une révision de
    silicium différente sur la carte d'évaluation (registre absent) ; perte du debug liée au cavalier de
    boot et à l'alimentation USB.
  \item Le pont diviseur de la CTN étant alimenté par la sortie de chauffe commutée, le nœud ADC est un
    signal carré à 100\,Hz pendant la PWM : échantillonnage synchronisé sur le temps haut (ISR à
    mi-impulsion, moyenne sur 16 mesures), qui a ramené la dispersion dans les limites du besoin de régulation.
\end{itemize}
\cvenv{Environnement\,:}{carte TLD4020-3STD\_EVAL, IAR Embedded Workbench, Infineon MCU Configuration Wizard,
  J-Link embarqué, printf sur LIN (MobaXterm, 19\,200 bauds 8N1)}}

\newcommand{\BlkForviaMeli}{%
\cvsub{Portage du driver LED MeLiBu : Melexis MLX81116 vers MLX81143}
\begin{itemize}
  \item Portage d'un driver MeLiBu~1.x existant sur RP2040 vers le MLX81143 : encodage de sous-adresse en
    mode étendu (14 LED par nœud), encodage PWM linéaire 15 bits vérifié sur datasheet et sur matériel ;
    panneau 22 × 7 (154 LED RVB).
  \item Correction des défauts de timing bus : génération du break en attente active sous interruptions
    masquées, section critique réduite au seul break pour ne plus perdre d'octets en réception UART,
    vidage de la FIFO TX avant chaque trame pour éviter les collisions ; correction de la broche de retour bus.
  \HW{\item Correction du câblage désordonné des LED dans la couche de mapping matériel, séparée de la
    convention d'entrée ; récepteur PIO utilisé pour confirmer les acquittements des esclaves.}
\end{itemize}
\cvenv{Environnement\,:}{RP2040, Raspberry Pi Pico SDK, machines d'états PIO, C, MeLiBu~1.x}}

\newcommand{\BlkForviaLight}{%
\cvsub{Démonstrateur d'éclairage d'ambiance intérieur (responsable logiciel)}
\begin{itemize}
  \item Extension du logiciel de pilotage d'un buck de démonstration intérieur automobile, de quelques
    afficheurs LED à neuf gateways UART STM32 pilotant des rubans WS2812 répartis en plusieurs zones.
  \item Application Python/PyQt5 multithread : décodage des animations MP4 avec OpenCV, correspondance entre
    identifiants LED et pixels de la frame, envoi de chaque frame par un thread par gateway à 921\,600 bauds,
    synchronisés par une barrière ; trames série propriétaires (ASCII avec compression delta).
  \item Application tablette Android (Kotlin, Jetpack Compose) qui pilote zones, couleurs et presets de
    luminosité en Bluetooth série, via un protocole de commandes JSON avec codes d'erreur ; mode de secours au clavier.
  \item Adaptation minimale du firmware des gateways STM32 (filtrage par décalage d'ID si deux gateways partagent une UART).
\end{itemize}
\cvenv{Environnement\,:}{Python 3, PyQt5, OpenCV, Kotlin/Jetpack Compose, STM32 (C), WS2812, UART, Bluetooth SPP, Git}}

% ---------------------------------------------------------------- KMUTT
\newcommand{\BlkKmutt}{%
\cventry{Stagiaire systèmes embarqués}{KMUTT --- MDT Lab}{Bangkok, Thaïlande}{avr.\ -- août 2025}
\begin{itemize}
  \item Traitement d'image pour un projet d'agriculture de précision : classification de fruits par transfer
    learning sur VGG16 (Keras/TensorFlow), 98,4\,\% de bonnes classifications sur un jeu de test de 4\,599 images
    (Fruits-360, 131 classes), et chaîne de détection OpenCV par contours et couleurs.
  \HW{\item Montée en compétence FPGA : conception d'un microcontrôleur 8 bits en VHDL (unité de contrôle
    fetch--decode--execute à 12 états, ALU avec drapeaux NZVC, ROM/RAM/E-S mappées en mémoire) d'après
    l'architecture de référence de LaMeres ; simulé sous Vivado 2024.2 / XSim.}
\end{itemize}
\cvenv{Environnement\,:}{Python, TensorFlow/Keras, OpenCV\HW{, VHDL, Vivado 2024.2}}}

% ---------------------------------------------------------------- projets académiques
\newcommand{\BlkAcadTherm}{%
\cvproject{Thermomètre temps réel : PIC24 bare-metal, puis FreeRTOS sur PIC32MZ}{Polytech, seul puis en binôme}{2025 -- 2026}{github.com/teo1747/pic24-32-rt-thermometer}
\begin{itemize}
  \item V1 (seul) : PIC24FJ1024GB610 en bare-metal, horodatage RTCC et journalisation en EEPROM SPI.
    V2 (binôme) : chaîne FreeRTOS de tâches acquisition, analyse et affichage sur PIC32MZ2048EFH100, reliée par
    des queues et des sémaphores binaires ; messages horodatés ; bascule °C/°F par bouton.
\end{itemize}
\cvenv{Environnement\,:}{Explorer 16/32, MPLAB X, XC16/XC32, Harmony v3/MCC, FreeRTOS 10.5, capteur TC1047A, débogueur embarqué}}

\newcommand{\BlkAcadZephyr}{%
\cvproject{Étude de faisabilité Zephyr RTOS sur STM32}{TP Polytech}{2025 -- 2026}{github.com/teo1747/zephyr-stm32-sensor-hub}
\begin{itemize}
  \item Trois threads se partagent un bus I²C : BMP280 (température/pression), MPU6050 (centrale inertielle
    6 axes) et OLED SSD1306 via l'API display et CFB ; données partagées protégées par mutex ; carte
    décrite en overlay devicetree (dépôt reconstruit d'après le rapport).
\end{itemize}
\cvenv{Environnement\,:}{NUCLEO-G474RE, Zephyr, West, CMake, Ninja, STM32CubeProgrammer, console série}}

\newcommand{\BlkAcadSram}{%
\cvproject{Contrôleur de SRAM ZBT en VHDL}{cours SoC Polytech, binôme}{2025 -- 2026}{github.com/teo1747/vhdl-zbt-sram-controller}
\begin{itemize}
  \item Contrôleur à machine d'états pour une SRAM ZBT Micron 512K × 36 : lectures/écritures simples
    et en rafale de 4 mots sur un bus bidirectionnel de 36 bits (IOBUF Xilinx) ; banc de test
    auto-vérifiant contre le modèle Micron.
\end{itemize}
\cvenv{Environnement\,:}{VHDL, Vivado, cible Zynq-7000, modèle de simulation Micron}}

\newcommand{\BlkAcadVideo}{%
\cvproject{Chaîne vidéo VGA temps réel sur ZedBoard}{Polytech, binôme}{2025 -- 2026}{github.com/teo1747/zedboard-vga-video-pipeline}
\begin{itemize}
  \item Chaîne 640 × 480 assemblée sous IP integrator : Clocking Wizard, Video Timing Controller, IP vidéo
    AXI4-Stream, inverseur VHDL et IP d'incrustation Vitis HLS. Synchros vérifiées sur carte à l'ILA ;
    la démo finale n'affichait pas d'image faute de maître AXI-Lite pour démarrer les IP HLS, depuis
    rendues autonomes (nouveau test sur carte à faire).
\end{itemize}
\cvenv{Environnement\,:}{Zynq-7000 (xc7z020), Vivado/Vitis HLS 2022.2, ILA, contraintes XDC}}

\newcommand{\BlkAcadFilter}{%
\cvproject{Filtre d'image 3×3 sur FPGA}{Polytech, seul}{nov.\ 2024 -- janv.\ 2025}{github.com/teo1747/vhdl-2d-image-filter}
\begin{itemize}
  \item Filtres moyenneur et gaussien (1-2-1/2-4-2/1-2-1) en flux sur pixels 8 bits : line buffers sur
    l'IP FIFO Generator, arbre d'additionneurs pipeliné avec croissance en bits explicite
    (9, 10 puis 11 bits) ; simulation complète sur l'image Lena, sorties comparées à une référence logicielle.
\end{itemize}
\cvenv{Environnement\,:}{VHDL, Vivado 2024.1, Artix-7 (xc7a100t), XSim, MATLAB}}

\newcommand{\BlkAcadTemp}{%
\cvproject{Moniteur de température sur FPGA}{Polytech, binôme}{2024 -- 2025}{github.com/teo1747/nexys4ddr-temperature-monitor}
\begin{itemize}
  \item Intégration de cœurs I²C (400\,kHz) et UART (115\,200 bauds) pour lire l'ADT7420 de la carte et l'afficher
    sur les afficheurs 7 segments ; timing tenu à 100\,MHz (WNS +5,4\,ns, 1\,936 LUT) sur Nexys 4 DDR.
\end{itemize}}

\newcommand{\BlkAcadMifare}{%
\cvproject{Porte-monnaie MIFARE Classic}{Polytech, seul}{oct.\ -- nov.\ 2024}{github.com/teo1747/mifare-classic-wallet-qt}
\begin{itemize}
  \item Application Qt/C++ : lecture/écriture de l'identité sur la carte et solde en bloc valeur avec bloc
    de sauvegarde, séparation des droits clé A/clé B ; testée avec un lecteur USB ODALID sur cartes physiques.
\end{itemize}}

\newcommand{\BlkAcadRpi}{%
\cvproject{Périphériques sous Linux embarqué sur Raspberry Pi}{Polytech, équipe}{mars 2025}{}
\begin{itemize}
  \item Drivers C en espace utilisateur pour modules Pmod : capteur de lumière SPI via spidev (mode 3, 2\,MHz),
    capteur d'humidité I²C via i2c-dev, GPIO via sysfs ; compilation par Make.
\end{itemize}}

% ---------------------------------------------------------------- projets personnels
\newcommand{\BlkPersonalNote}{%
\vspace{1pt}{\fontsize{8.4pt}{10.6pt}\selectfont\color{muted}\FP{Projets personnels. Je fixe l'architecture et
le plan de vérification ; l'implémentation est assistée par des outils de programmation IA. Vérifiés en
simulation et en émulation, pas encore sur cartes physiques.}\NF{En cours. Je définis l'architecture, les
spécifications et la stratégie de vérification ; l'implémentation est assistée par des outils de
programmation IA. Validation à ce stade en émulation (QEMU, EmbSIM), pas encore sur cartes physiques.}\par}}

\newcommand{\BlkRtos}{%
\cvproject{EmbLinkRTOS --- RTOS préemptif pour microcontrôleurs}{personnel}{depuis oct.\ 2026}{github.com/EmbLinkOs/EmbLinkRTOS}
\begin{itemize}
  \item Ordonnanceur préemptif à priorités fixes (bitmap de listes prêtes, FIFO à priorité égale) ;
    port Cortex-M (commutation PendSV/PSP, sauvegarde FPU paresseuse, sections critiques BASEPRI),
    ports RISC-V et AVR.
  \item Sémaphores, mutex avec héritage et plafond de priorité, timeouts ; tests de conformité et
    différentiels contre un modèle Python de l'ordonnanceur.
  \item Build CMake/Kconfig avec définitions de cartes (NUCLEO-F446RE, Arduino Uno, ESP32-C3) ; CI GitHub
    Actions avec ASan/UBSan/TSan, exécutions QEMU et EmbSIM.
\end{itemize}}

\newcommand{\BlkEmbLinkOS}{%
\cvproject{EmbLinkOS --- système d'exploitation x86-64 écrit de zéro}{personnel}{depuis mai 2026}{github.com/EmbLinkOs/EmbLinkOs}
\begin{itemize}
  \item Bootloader BIOS en deux étages en NASM (mode réel, protégé, puis long) et chargeur UEFI PE32+ ;
    noyau higher-half en C : pagination 4 niveaux, APIC/IO-APIC, SMP, processus ring 3, appels système, chargeur ELF.
  \item Système de fichiers B-tree copy-on-write (CRC32C, snapshots, AES-XTS) vérifié contre un oracle
    Python ; FAT32, VFS ; drivers PCI, AHCI, NVMe et USB ; portages de newlib et TCC.
  \item Autotests noyau, tests unitaires hôte sous ASan/UBSan, tests d'intégration pilotés sous QEMU ; debug GDB distant.
\end{itemize}}

\newcommand{\BlkEmbCC}{%
\cvproject{EmbCC --- chaîne de compilation croisée C}{personnel}{depuis juil.\ 2026}{github.com/EmbLinkOs/EmbCC}
\begin{itemize}
  \item Compilateur C99, assembleur et éditeur de liens avec back-ends Cortex-M, RISC-V, AVR et x86-64 ;
    utilitaires firmware (génération d'en-têtes/startup depuis SVD, empaquetage HEX/UF2).
  \item Vérification : encodages comparés à llvm-mc/GNU as, code généré exécuté sous QEMU et comparé à
    GCC/Clang ; CI GitHub Actions sous Linux et macOS avec GCC et Clang.
\end{itemize}}

\newcommand{\BlkEmbSim}{%
\cvproject{EmbSIM --- simulateur de microcontrôleurs et de circuits}{personnel}{depuis oct.\ 2026}{github.com/EmbLinkOs/EmbSIM}
\begin{itemize}
  \item Simulateurs de jeux d'instructions et modèles de périphériques (Cortex-M, AVR ; GPIO, timers, USART, SPI,
    I²C, ADC, DMA des STM32F1/F4) couplés à un solveur de circuit par analyse nodale ; la CI exécute des
    exemples STM32Cube, Arduino et Pico SDK non modifiés. Stade précoce.
\end{itemize}}

\newcommand{\BlkAccess}{%
\cvproject{Terminal de contrôle d'accès RFID/PIN sous FreeRTOS}{personnel}{oct.\ 2026}{github.com/teo1747/stm32-freertos-access-control}
\begin{itemize}
  \item Drivers au niveau registre (sans HAL constructeur) sur STM32F446RE : RFID MFRC522 en SPI, SSD1306, RTC
    DS3231 et EEPROM 24LC256 en I²C, clavier, servo en PWM, watchdog ; six tâches FreeRTOS avec queues, mutex,
    event groups et stream buffers ; PIN hachés SHA-256 salés et étirés ; journal protégé par CRC.
  \item Tests unitaires hôte avec sanitizers et scénarios système sous EmbSIM ; pas encore testé sur carte physique.
\end{itemize}
\cvenv{Environnement\,:}{NUCLEO-F446RE, FreeRTOS 11, CMSIS, arm-none-eabi-gcc, CMake, Ninja}}

\newcommand{\BlkAccel}{%
\cvproject{Accélérateur de filtrage d'image pour Zynq-7000 avec DMA AXI}{personnel}{oct.\ 2026}{github.com/teo1747/zynq-image-filter-accelerator}
\begin{itemize}
  \item IP de filtrage 3×3 en flux en VHDL (moyenneur, gaussien, Sobel), un pixel par cycle : line buffers,
    pipeline à 4 étages, back-pressure AXI4-Stream complète, tramage TLAST/TUSER ; registres de contrôle et
    d'état AXI4-Lite ; IP packagée pour IP integrator.
  \item Block design PS7 Zynq + DMA AXI scripté en TCL jusqu'au bitstream et au XSA ; application Vitis
    bare-metal qui envoie les images depuis la DDR, vérifie chaque pixel contre le modèle C et compare les
    temps matériel et logiciel.
  \item Bit-exact face à une référence C dans un banc de test GHDL auto-vérifiant avec blocages aléatoires
    sur les deux flux (14 images, exécuté en CI). Essai sur ZedBoard à venir.
\end{itemize}
\cvenv{Environnement\,:}{VHDL, C, TCL, Vivado, Vitis, GHDL, Zynq-7000 (xc7z020)}}

% ---------------------------------------------------------------- formation & langues
\newcommand{\BlkEducation}{%
\section{Formation \& langues}
\cvedu{Diplôme d'ingénieur (Bac+5) --- Systèmes embarqués}{Polytech Dijon}{Dijon}{2023 -- sept.\ 2026}
{\color{muted}Diplômé. Dernière année en alternance chez FORVIA. \FP{Enseignements : conception numérique, SoC sur Zynq (Vivado, Vitis, Vitis HLS), temps réel.}\NF{Enseignements : électronique, PCB, temps réel, SoC/FPGA.}\par}
\cvedu{Classes préparatoires (2 ans) --- mathématiques, physique, électronique}{IUC}{\iucplace}{\iucdates}
\vspace{4pt}
\cvskill{Langues}{Anglais : langue maternelle (anglophone), TOEIC C1 (2024)\sep Français : courant (B2)}
\cvskill{Certification}{Cisco CCNA 1 (juillet 2024)}}
\newcommand{\iucplace}{Cameroun}
\newcommand{\iucdates}{2021 -- 2023}

\newcommand{\BlkLanguages}{} % intégré à \BlkEducation
